\documentclass[11pt]{article}
\usepackage[T1]{fontenc}
\usepackage{lmodern}
\usepackage[margin=1in]{geometry}
\usepackage{amsmath,amssymb,amsthm,mathtools}
\usepackage{microtype}
\usepackage[colorlinks=true,linkcolor=blue,citecolor=blue,urlcolor=blue]{hyperref}
\numberwithin{equation}{section}
\newtheorem{theorem}{Theorem}[section]
\newtheorem{lemma}[theorem]{Lemma}
\newtheorem{proposition}[theorem]{Proposition}
\newcommand{\eps}{\varepsilon}
\newcommand{\sech}{\operatorname{sech}}
\newcommand{\Span}{\operatorname{span}}
\newcommand{\HH}{\mathcal H^2}
\newcommand{\LL}{\mathcal L}
\newcommand{\norm}[1]{\left\|#1\right\|}
\newcommand{\ip}[2]{\left\langle #1,#2\right\rangle}
\title{Exact uniform control time for a centered viscous Burgers shock}
\author{Kaiwen Jin\thanks{Email: \href{mailto:yui.ui.siki@gmail.com}{yui.ui.siki@gmail.com}.}}
\date{October 2026}

\begin{document}
\maketitle
\begin{abstract}
We determine the minimal time for boundary null control of the Burgers
equation linearized at the stationary shock
$-\tanh(x/(2\eps))$ on $(-L,L)$, with a control at the left endpoint.
For arbitrary $L^2$ initial data, the worst-case $L^2$ control cost remains
bounded as the viscosity tends to zero at every time $T>4\sqrt2L$ and
diverges at every time $0<T<4\sqrt2L$.
Thus the exact uniform control time is $4\sqrt2L$.
The upper bound separates the near-zero shock mode from the fast spectral
modes and estimates their coefficient norms through the exact spectral
determinant. The lower bound uses an inverse exponential Gram matrix:
its alternating coefficients combine with the eigenfunctions' endpoint
signs to prevent cancellation in a common boundary layer.
\end{abstract}

\section{Introduction}\label{sec:intro}

Consider disturbances around the centered stationary viscous Burgers shock
\[
 U^\eps(x)=-\tanh\!\left(\frac{x}{2\eps}\right),\qquad -L<x<L,
\]
where $L>0$ and $\eps>0$. Their linearized equation with one boundary
actuator is
\begin{equation}\label{eq:primal}
 \begin{cases}
  u_t-\eps u_{xx}+\partial_x(U^\eps u)=0,&0<t<T,\ -L<x<L,\\
  u(t,-L)=h(t),\qquad u(t,L)=0,&0<t<T,\\
  u(0,x)=u_0(x),&-L<x<L.
 \end{cases}
\end{equation}
We use the usual transposition solution for real $u_0\in L^2(-L,L)$ and
real $h\in L^2(0,T)$; Appendix~\ref{sec:duality} recalls its construction and
the exact control--observation duality. No pointwise compatibility between
the initial datum and the boundary values is imposed.

The question is how long the actuator must act to drive every disturbance
to zero with a cost that stays bounded in the vanishing-viscosity limit.
For fixed $L$, define
\begin{align}
 C(T,\eps)
 &=\sup_{\norm{u_0}_{L^2(-L,L)}=1}
   \inf\{\norm h_{L^2(0,T)}:u(T)=0\},\label{eq:cost}\\
 T_{\rm unif}
 &=\inf\{T>0:\limsup_{\eps\downarrow0}C(T,\eps)<\infty\}.
 \label{eq:time}
\end{align}
The infimum over an empty control set is $+\infty$.
Although the limiting transport picture suggests a geometric time scale,
the small-viscosity spectral structure also enters this uniform problem.
Laheurte~\cite[Theorem 1.1 and Remark 1.2]{Laheurte} obtained
\begin{equation}\label{eq:old}
 (4\sqrt2-2)L\le T_{\rm unif}\le4\sqrt3L
\end{equation}
for the same centered shock and boundary conditions.
The present result identifies the infimum in \eqref{eq:time} without
restricting the initial data or changing the boundary control norm.

\begin{theorem}\label{thm:main}
For every $L>0$, the problem \eqref{eq:primal} satisfies
\[
 T_{\rm unif}=4\sqrt2L.
\]
More precisely:
\begin{enumerate}
 \item For every $T>4\sqrt2L$, there are $\eps_0>0$ and $M_{L,T}<\infty$
 such that, for every $0<\eps<\eps_0$ and every $u_0\in L^2(-L,L)$,
 there is a real null control $h\in L^2(0,T)$ with
 $\norm h_{L^2(0,T)}\le M_{L,T}\norm{u_0}_{L^2(-L,L)}$.
 \item For every $0<T<4\sqrt2L$, there are $\eps_0>0$ and $c_{L,T}>0$
 such that
 \begin{equation}\label{eq:lower-main}
  C(T,\eps)\ge c_{L,T}\eps^{13/2}
       \exp\!\left(\frac{4\sqrt2L-T}{4\eps}\right),
       \qquad 0<\eps<\eps_0.
 \end{equation}
\end{enumerate}
All constants depend only on the displayed fixed parameters.
\end{theorem}

The first assertion concerns every $L^2$ initial datum; in the second,
the adverse datum may depend on $\eps$, as the supremum in \eqref{eq:cost}
allows. These strict-side statements determine the exact infimum.
They leave the control cost at $T=4\sqrt2L$ open.

The main difficulty in the upper bound is that one eigenvalue approaches
zero while all remaining eigenvalues grow on the scale $\eps^{-1}$.
The slow mode therefore cannot be absorbed into an estimate for the fast
modes. We separate its boundary observation from the fast exponential
span by showing that bounded fast observations have weak limits supported
in a fixed finite time interval. Hardy-space projections make this
support statement precise. The exact spectral determinant then bounds
all fast coefficient functionals, including arbitrarily large indices,
with the exponent that gives the time $4\sqrt2L$.

For the lower bound, cancellation between fast modes is the essential
issue. A column of the inverse Gram matrix of their temporal exponentials
has alternating signs. Centered reflection symmetry gives the opposite
endpoint derivatives of the corresponding spatial eigenfunctions with
exactly the complementary signs. Their products consequently have one
sign in a common endpoint layer. A scalar separation product determines
the matching exponential cost.
The inverse-Gram and Hardy model-space approach is informed by
Koike and Laheurte's work on constant-drift transport--diffusion
\cite[Sections 3, 5 and 6.1]{KoikeLaheurte}.
Here the determinant and the endpoint signs are those of the centered
shock. The theorem concerns this linearized equation; it makes no claim
for nonlinear Burgers dynamics or off-center shocks.

Section~\ref{sec:spectral} derives the spectrum and boundary-normalized
eigenfunctions. Section~\ref{sec:upper} proves uniform observability above
the asserted time. Section~\ref{sec:lower} constructs the matching adverse
tests. Appendix~\ref{sec:duality} passes between observability and controls.
Throughout, $D=2L$ is the interval length. Spatial norms $\norm{\cdot}_2$
use Lebesgue measure on $(-L,L)$, and $\norm f_J$ denotes the $L^2$
norm on a time interval $J$. Complex inner products are linear in the
first argument; the original data and controls are real. Constants may
change from line to line but never depend on $\eps$, a spectral index,
or a spectral truncation unless this dependence is stated explicitly.
All small-viscosity estimates hold on an interval $(0,\eps_0)$.

\section{Spectrum and normalized boundary observations}\label{sec:spectral}

The adjoint realization and its domain are
\begin{equation}\label{eq:adjoint}
 A_\eps^*=-\eps\partial_x^2+
          \tanh\!\left(\frac{x}{2\eps}\right)\partial_x,\qquad
 X_\eps=D(A_\eps^*)=H^2(-L,L)\cap H_0^1(-L,L).
\end{equation}
For $z_t+A_\eps^*z=0$ with $z(0)=z_0\in X_\eps$, the observed flux is
\begin{equation}\label{eq:flux}
 q(t)=\eps z_x(t,-L).
\end{equation}
Integration by parts gives
\begin{equation}\label{eq:duality}
 \ip{u(T)}{z_0}
  =\ip{u_0}{z(T)}+\int_0^T h(t)\overline{q(T-t)}\,dt.
\end{equation}
Thus an estimate $\norm{z(T)}_2\le M\norm q_{(0,T)}$ provides a null
control of norm at most $M\norm{u_0}_2$, while every nonzero smooth
adjoint test gives
\begin{equation}\label{eq:ratio}
 C(T,\eps)\ge\frac{\norm{z(T)}_2}{\norm q_{(0,T)}}.
\end{equation}
The extension to all $L^2$ data is proved in Appendix~\ref{sec:duality}.

\subsection{A first-order reduction}

Set
\[
 b=\frac{L}{2\eps},\qquad a=\frac{\tanh b}{2\eps},\qquad
 c(x)=\cosh\!\left(\frac{x}{2\eps}\right).
\]
Direct conjugation gives the self-adjoint Dirichlet operator
\begin{equation}\label{eq:H}
 c^{-1}A_\eps^*c=H_\eps
 =-\eps\partial_x^2+\frac1{4\eps}
       -\frac1{2\eps}\sech^2\!\left(\frac{x}{2\eps}\right).
\end{equation}
With the formal first-order factors
\[
 B=\eps\partial_x+\tfrac12\tanh(x/(2\eps)),\qquad
 B^\dagger=-\eps\partial_x+\tfrac12\tanh(x/(2\eps)),
\]
one has $\eps H_\eps=B^\dagger B$ and
$BB^\dagger=-\eps^2\partial_x^2+1/4$.
In particular, $H_\eps$ is positive: a Dirichlet solution of $Bv=0$
would be a multiple of $\sech(x/(2\eps))$ and hence zero.

If $H_\eps v=\lambda v$ with zero endpoint values, $f=Bv$ solves
\begin{equation}\label{eq:robin}
 -\eps^2f''+\tfrac14f=\eps\lambda f,\qquad
 f'(-L)=-af(-L),\qquad f'(L)=af(L).
\end{equation}
Conversely, $B^\dagger f$ transfers every nonzero eigenvalue of
\eqref{eq:robin} back to the Dirichlet problem.
The Robin eigenfunction $f=c$ has eigenvalue zero and is annihilated by
$B^\dagger$; it is not an eigenfunction of $H_\eps$.
This reduction is the Darboux structure used in
\cite[Section 2]{Laheurte}; the needed identities are derived here.

For small enough $\eps$, there is exactly one Dirichlet eigenvalue below
$1/(4\eps)$. Denote it by $\lambda_0$. It has the form
\begin{equation}\label{eq:slow}
 \lambda_0=\frac1{4\eps}-\eps\kappa^2,\qquad
 \kappa\coth(\kappa L)=a,\qquad 0<\kappa<\frac1{2\eps}.
\end{equation}
Indeed, a general hyperbolic Robin solution is a combination of
$\cosh(\kappa x)$ and $\sinh(\kappa x)$. Adding and subtracting the
endpoint conditions gives respectively
$\kappa\tanh(\kappa L)=a$ and $\kappa\coth(\kappa L)=a$.
Both left sides are strictly increasing; for the odd equation the
derivative has the sign of $\sinh(2\kappa L)-2\kappa L>0$.
The even root is $1/(2\eps)$ and produces the annihilated mode.
The odd root exists when $a>1/L$, is unique, and is smaller than that
even root.

All remaining eigenvalues are
\begin{equation}\label{eq:fast}
 \lambda_k=\frac1{4\eps}+\eps\xi_k^2,\qquad
 D\xi_k-2\arctan(\xi_k/a)=k\pi,\qquad k\ge1.
\end{equation}
For $Da>2$ the phase on the left is strictly increasing from zero.
The zero frequency is excluded because the limiting characteristic
expression is $2a-a^2D\ne0$. Consequently
\begin{equation}\label{eq:brackets}
 \frac{k\pi}{D}<\xi_k<\frac{(k+1)\pi}{D},\qquad k\ge1.
\end{equation}

\subsection{Endpoint normalization and spatial bounds}

For a general complex parameter $\lambda$ put
\[
 f_\lambda(x)=\cos(\xi(x+L))-\frac a\xi\sin(\xi(x+L)),\qquad
 \xi^2=\frac\lambda\eps-\frac1{4\eps^2},
\]
and define
\begin{equation}\label{eq:S}
 S_\lambda(x)=
 \frac{-\eps f_\lambda'(x)
              +\frac12\tanh(x/(2\eps))f_\lambda(x)}{\lambda}.
\end{equation}
This is the entire fundamental solution of
$(H_\eps-\lambda)S_\lambda=0$ with
$S_\lambda(-L)=0$, $S_\lambda'(-L)=1$.
To see the removability in \eqref{eq:S}, cosine and $\sin(\xi y)/\xi$
are entire in $\xi^2$, and $f_0=c/c(L)$ is annihilated by $B^\dagger$.
At $-L$ the derivative of the numerator equals $\lambda$, using
$\eps a^2+(4\eps)^{-1}\sech^2 b=(4\eps)^{-1}$.
The initial-value characterization also gives uniqueness.

Normalize each adjoint eigenfunction by its left flux:
\begin{equation}\label{eq:chi}
 \chi_j(x)=\frac{c(x)}{\eps c(L)}S_{\lambda_j}(x),\qquad j\ge0.
\end{equation}
At a fast root in \eqref{eq:fast},
$f_{\lambda_k}(L)=(-1)^{k+1}$.
At either Dirichlet endpoint the derivative of
$B^\dagger f_{\lambda_k}$ equals $\lambda_k f_{\lambda_k}$.
Hence the normalization gives the exact identities
\begin{equation}\label{eq:endpoint}
 \eps\chi_j'(-L)=1\quad(j\ge0),\qquad
 \eps\chi_k'(L)=(-1)^{k+1}\quad(k\ge1).
\end{equation}

\begin{lemma}\label{lem:norms}
For all sufficiently small $\eps$, uniformly in $k\ge1$,
\begin{equation}\label{eq:norms}
 \norm{\chi_k}_2\le\frac{C_L}{k\sqrt\eps},\qquad
 \norm{\chi_0}_2\le2\sqrt D,\qquad \lambda_0\longrightarrow0.
\end{equation}
\end{lemma}
\begin{proof}
For fast modes, \eqref{eq:S} and $|\tanh(x/(2\eps))|\le\tanh b$ yield
\[
 |\chi_k(x)|\le\frac{c(x)}{c(L)}
       \frac{(\xi_k+a)^2}{\xi_k\lambda_k}
 \le\frac{2}{\eps\xi_k}\frac{c(x)}{c(L)}.
\]
The last inequality uses $a\le1/(2\eps)$ and
$(s+1/2)^2\le2(s^2+1/4)$.
For small $\eps$,
$\int_{-L}^L(c(x)/c(L))^2\,dx\le8\eps$.
Together with \eqref{eq:brackets}, this proves the first bound.

For the slow mode, set $\gamma=2\eps\kappa$,
$\beta=(2\eps)^{-1}-\kappa$ and $A=(2\eps)^{-1}+\kappa$.
Equation~\eqref{eq:slow} becomes
$\gamma=\tanh b\,\tanh(\gamma b)$.
Its positive root satisfies $\gamma\ge1/2$ for large $b$, and
\[
 1-\gamma\le(1-\tanh b)+(1-\tanh(\gamma b))\le4e^{-b}.
\]
Thus $\lambda_0\le2\eps^{-1}e^{-b}\to0$ and
$\beta\le2\eps^{-1}e^{-b}\to0$.
The hyperbolic product identities and \eqref{eq:slow} give the positive
adjoint ground function
\[
 w(x)=\frac{1+\gamma}{4}
 \left(\cosh(\beta x)
       -\frac{\cosh(\beta L)}{\cosh(AL)}\cosh(Ax)\right).
\]
Positivity in the interior follows because
$\cosh(Ax)/\cosh(\beta x)$ increases for $0\le x\le L$.
Normalizing $w$ by its left flux, and differentiating there, gives
\[
 0\le\chi_0(x)=\frac{w(x)}{\eps w'(-L)}
 \le\frac1{\eps[A\tanh(AL)-\beta\tanh(\beta L)]}\le2
\]
for small $\eps$; the denominator tends to one.
This proves the remaining statements.
\end{proof}

\subsection{Spectral completeness and the derivative trace}\label{sec:density}

Regular separated Sturm--Liouville theory applies to \eqref{eq:H}:
the interval is finite, the leading coefficient $\eps$ is positive,
the potential is real and smooth, and the boundary conditions are
Dirichlet. Its self-adjoint realization has compact resolvent and a
complete orthonormal eigenbasis $(\phi_j)_{j\ge0}$
\cite[Theorems 9.6 and 9.10]{Teschl}.
Its domain is precisely $H^2\cap H_0^1$.
For fixed $\eps$, if $V$ is the potential in \eqref{eq:H},
\[
 \norm{v''}_2\le\eps^{-1}(\norm{H_\eps v}_2+\norm V_\infty\norm v_2),
 \qquad \norm{v'}_2^2\le\norm{v''}_2\norm v_2.
\]
These estimates, with the reverse bound from the bounded coefficients,
make the graph norm equivalent to the $H^2$ norm.

Writing $v=\sum a_j\phi_j$, the domain condition is
$\sum(1+\lambda_j^2)|a_j|^2<\infty$.
One direction follows by taking the coefficients of $H_\eps v$;
the other follows by closedness from finite spectral sums.
The expansion $e^{-tH_\eps}v=\sum e^{-\lambda_jt}a_j\phi_j$ then
converges in the graph norm uniformly for $0\le t\le T$, since the
tail is bounded by its value at $t=0$.
Multiplication by $c$ and $c^{-1}$ preserves $H^2\cap H_0^1$ at each
fixed $\eps$. Thus finite sums of the $\chi_j$ approximate every
adjoint initial test in $X_\eps$, and their solutions converge in
$C([0,T];H^2)$.

For $g\in H^1(-L,L)$, its absolutely continuous representative satisfies
\[
 |g(-L)|\le D^{-1/2}\norm g_2+\sqrt D\norm{g'}_2.
\]
This follows by averaging
$|g(-L)|\le|g(x)|+\int_{-L}^x|g'(r)|\,dr$.
Applying it to $g=v'$ proves continuity of the derivative trace on
$H^2$. The approximating observed fluxes therefore converge in
$L^2(0,T)$.
Constants in this density argument may depend on the fixed $\eps$;
we apply it only after proving estimates uniform in the spectral
truncation and in small viscosity.

\section{Uniform observability above the critical time}\label{sec:upper}

\subsection{The determinant and spectral separation products}

The coefficient estimates for a fast exponential sum depend on its
spectral separation. We obtain an all-index bound from the exact
Dirichlet determinant
\begin{equation}\label{eq:Delta}
 \Delta_\eps(\lambda)=S_\lambda(L)
 =\frac\eps\lambda
  \left[\left(\xi-\frac{a^2}{\xi}\right)\sin(D\xi)
                  +2a\cos(D\xi)\right].
\end{equation}
It is entire, and the explicit formula bounds its growth by a polynomial
times $\exp(C_\eps\sqrt{|\lambda|})$, so its order is at most $1/2$.
Its zeros are exactly the positive Dirichlet eigenvalues.
Differentiating the differential equation in $\lambda$ and integrating
the Wronskian gives
\[
 \eps S_{\lambda_j}'(L)\Delta_\eps'(\lambda_j)
       =\int_{-L}^L S_{\lambda_j}(x)^2\,dx>0.
\]
Here $S_{\lambda_j}'(L)\ne0$ by initial-value uniqueness, so every zero
is simple. At zero the fundamental solution is explicitly
\[
 S_0(x)=\frac{\sech(x/(2\eps))}{c(L)}
                      \int_{-L}^x c(r)^2\,dr,\qquad
 \Delta_\eps(0)=c(L)^{-2}\int_{-L}^L c(r)^2\,dr>0.
\]
Hadamard factorization for an entire function of order below one
\cite[Section 6.4.3]{Weikard} consequently gives
\begin{equation}\label{eq:Hadamard}
 \Delta_\eps(\lambda)=\Delta_\eps(0)
                     \prod_{j\ge0}(1-\lambda/\lambda_j).
\end{equation}
There is no nonconstant exponential factor.

For $k\ge1$ define the full and fast separation products
\[
 \widetilde d_k=\prod_{\substack{j\ge0\\j\ne k}}
       \frac{|\lambda_j-\lambda_k|}{\lambda_j+\lambda_k},
 \qquad
 d_k=\prod_{\substack{j\ge1\\j\ne k}}
       \frac{|\lambda_j-\lambda_k|}{\lambda_j+\lambda_k}.
\]
They converge to positive numbers because
$\lambda_j\sim\eps\pi^2j^2/D^2$.
The analytic product criterion used here and below is
\cite[Section 1.10(ix)]{DLMF}: summable uniform deviations of the
factors on compact sets give locally uniform convergence, with no
additional zeros.
Equation~\eqref{eq:Hadamard} implies
\begin{equation}\label{eq:dk-full}
 d_k\ge\widetilde d_k
     =\frac{2\lambda_k|\Delta_\eps'(\lambda_k)|}
                         {\Delta_\eps(-\lambda_k)}.
\end{equation}

\begin{lemma}\label{lem:dk}
For all sufficiently small $\eps$ and every $k\ge1$,
\begin{equation}\label{eq:dk}
 d_k^{-1}\le C_L\eps e^{D\eta_k},\qquad
 \eta_k=\sqrt{\xi_k^2+\frac1{2\eps^2}}.
\end{equation}
\end{lemma}
\begin{proof}
The identity
\[
 \left(\xi-\frac{a^2}{\xi}\right)\sin(D\xi)+2a\cos(D\xi)
 =-\frac{\xi^2+a^2}{\xi}
                 \sin(D\xi-2\arctan(\xi/a))
\]
and \eqref{eq:fast} yield
\begin{align}
 2\lambda_k|\Delta_\eps'(\lambda_k)|
  &=\frac{D(\xi_k^2+a^2)-2a}{\xi_k^2},
       \label{eq:Delta-derivative}\\
 \Delta_\eps(-\lambda_k)
  &=\frac\eps{\lambda_k}
    \left[\left(\eta_k+\frac{a^2}{\eta_k}\right)\sinh(D\eta_k)
                         -2a\cosh(D\eta_k)\right].
       \label{eq:Delta-negative}
\end{align}
If $Da\ge4$, the first expression is at least
$D(\xi_k^2+a^2)/(2\xi_k^2)$.
The bracket in the second is positive by \eqref{eq:Hadamard}
and at most $(\eta_k+a^2/\eta_k)e^{D\eta_k}/2$.
Finally,
\[
 \frac{a^2}{\eta_k^2}\le\frac12,\qquad
 \frac{\eps\eta_k}{\lambda_k}
   =\eps\frac{\sqrt{\eps\lambda_k+1/4}}{\eps\lambda_k}
   \le2\sqrt2\,\eps.
\]
Combining these bounds with \eqref{eq:dk-full} proves \eqref{eq:dk},
uniformly in the index.
\end{proof}

\subsection{Exponential Gram matrices and the fast terminal estimate}

For $N\ge1$ let
\[
 G_{mn}=\int_0^\infty e^{-(\lambda_m+\lambda_n)t}\,dt
       =\frac1{\lambda_m+\lambda_n},\qquad 1\le m,n\le N,
\]
and let $d_{k,N}$ be the product defining $d_k$ restricted to
$1\le j\le N$, $j\ne k$.
The exponentials are linearly independent: a vanishing combination,
being continuous, vanishes everywhere, and its first $N$ derivatives
at zero give a Vandermonde system. Thus $G$ is positive definite.

The signed Cauchy inverse is
\begin{equation}\label{eq:Gram}
 (G^{-1})_{mn}
  =\frac{4(-1)^{m+n}\lambda_m\lambda_n}
   {(\lambda_m+\lambda_n)d_{m,N}d_{n,N}},\qquad
 (G^{-1})_{kk}=\frac{2\lambda_k}{d_{k,N}^2}.
\end{equation}
For a direct verification, put
$p_i=2\lambda_i\prod_{j\ne i}(\lambda_i+\lambda_j)/
(\lambda_i-\lambda_j)$.
Partial fractions give
\[
 Q(s)=\prod_{j=1}^N\frac{s-\lambda_j}{s+\lambda_j}
       =1-\sum_{j=1}^N\frac{p_j}{s+\lambda_j}.
\]
Its values at the roots and $Q'(\lambda_i)=1/p_i$ imply
\[
 G\operatorname{diag}(p)G=\operatorname{diag}(1/p),\qquad
 G^{-1}=\operatorname{diag}(p)G\operatorname{diag}(p).
\]
This gives both magnitudes and signs in \eqref{eq:Gram}.

If $f(t)=\sum_{k=1}^N c_ke^{-\lambda_kt}$, write
$\boldsymbol c=(c_1,\ldots,c_N)^T$.
Its norm squared is $\boldsymbol c^*G\boldsymbol c$.
The coefficient functional has squared norm $(G^{-1})_{kk}$.
Since $d_{k,N}\ge d_k$, we obtain
\[
 |c_k|\le\frac{\sqrt{2\lambda_k}}{d_k}\norm f_{(0,\infty)}.
\]
For the corresponding fast adjoint solution
$z_{\rm f}(t)=\sum c_ke^{-\lambda_kt}\chi_k$, this implies
\begin{equation}\label{eq:fast-terminal}
 \norm{z_{\rm f}(T)}_2\le
 \left(\sum_{k\ge1}e^{-\lambda_kT}
              \frac{\sqrt{2\lambda_k}\norm{\chi_k}_2}{d_k}\right)
                     \norm f_{(0,\infty)}.
\end{equation}
The sum over all indices bounds every finite truncation.
Concavity of the square root gives
\begin{equation}\label{eq:exponent}
 D\eta_k-\lambda_kT
 \le-\frac{T-2\sqrt2D}{4\eps}
       -\eps(T-D/\sqrt2)\xi_k^2.
\end{equation}
Indeed, $\eta_k\le1/(\sqrt2\eps)+(\eps/\sqrt2)\xi_k^2$.
Combining \eqref{eq:brackets}, \eqref{eq:norms}, \eqref{eq:dk}
and \eqref{eq:exponent}, for $T>2\sqrt2D$ we find
\begin{equation}\label{eq:fast-upper}
 \norm{z_{\rm f}(T)}_2\le
 C_{L,T}\eps^{-1/2}
    e^{-(T-2\sqrt2D)/(4\eps)}\norm f_{(0,\infty)}.
\end{equation}
For clarity, the prefactor at index $k$ is bounded by
$C_L(1/k+\eps)$, since
$\sqrt{\eps\lambda_k}/k\le C_L(1/k+\eps)$.
For $\eps\le1$ this prefactor is bounded by a constant, and the
remaining Gaussian sum $\sum_{k\ge1}e^{-c_{L,T}\eps k^2}$
is at most $C_{L,T}\eps^{-1/2}$.

\subsection{Finite-time comparison and a limit of fast observations}
\label{sec:Hardy}

We now compare the half-line norm in \eqref{eq:fast-upper} with the
measured norm on $(0,T)$, and identify the weak limits needed to separate
the slow mode. Let $\HH$ be the Hardy space on $\Re s>0$ with norm
\[
 \norm F_{\HH}^2=\sup_{\sigma>0}\frac1{2\pi}
                         \int_{\mathbb R}|F(\sigma+iy)|^2\,dy.
\]
The Laplace transform
$\LL f(s)=\int_0^\infty e^{-st}f(t)\,dt$
is unitary from $L^2(0,\infty)$ onto $\HH$.
Boundary values are understood in $L^2(dy/(2\pi))$, and the inverse
has the Fourier sign $+iyt$ and factor $1/(2\pi)$
\cite[Propositions 3.4--3.5]{KoikeLaheurte}.

An inner function $\Theta$ is bounded analytic on the right half-plane
with unit boundary modulus almost everywhere.
Its multiplier is an isometry, so $\Theta\HH$ is closed.
The model space $K_\Theta=\HH\ominus\Theta\HH$ has orthogonal projection
\begin{equation}\label{eq:projection}
 P_\Theta F=F-\Theta P_+(\overline\Theta F),\qquad F\in\HH,
\end{equation}
where $P_+$ projects boundary $L^2$ onto $\HH$.
This follows by taking the multiplier's adjoint:
$M_\Theta^*F=P_+(\overline\Theta F)$.
The evaluation kernel is $1/(s+\lambda)=\LL(e^{-\lambda t})$,
because $\ip F{1/(s+\lambda)}=F(\lambda)$.
When $\Theta(\lambda)=0$ this kernel belongs to $K_\Theta$.

For the finite fast span use the rational inner function
\[
 \Theta_{\eps,N}(s)=\prod_{k=1}^N\frac{\lambda_k-s}{\lambda_k+s}.
\]
Every $\LL f$ for a finite fast observation with indices at most $N$
belongs to its model space.
The constant-one Bernstein inequality
\cite[Theorem 5.2]{BaranovEtAl} states that, for an inner function
with bounded boundary derivative and $F\in K_\Theta$,
\[
 \norm{F'}_{\HH}\le\norm{\Theta'}_\infty\norm F_{\HH}.
\]
The cited theorem is stated on the upper half-plane; the rotation
$s=-iz$ has derivative modulus one. The boundary normalization multiplies
both norms by the same factor and hence leaves the constant unchanged.
Our finite products satisfy its hypotheses, and on the boundary
\[
 |\Theta_{\eps,N}'(iy)|
    =2\sum_{k=1}^N\frac{\lambda_k}{\lambda_k^2+y^2}
    \le2\sum_{k\ge1}\lambda_k^{-1}<2D.
\]
Here the spectral brackets and the decreasing integral comparison give
\begin{equation}\label{eq:sum}
 \sum_{k\ge1}\lambda_k^{-1}
 <\int_0^\infty
     \frac{dr}{1/(4\eps)+\eps\pi^2r^2/D^2}=D.
\end{equation}
As $-\LL^{-1}F'=tf$, Bernstein yields
$\norm{tf}_{(0,\infty)}\le2D\norm f_{(0,\infty)}$.
The tail beyond $T$ is therefore at most
$4D^2T^{-2}\norm f_{(0,\infty)}^2$.
For every $T>2D$,
\begin{equation}\label{eq:comparison}
 \norm f_{(0,\infty)}^2
    \le\frac{T^2}{T^2-4D^2}\norm f_{(0,T)}^2
\end{equation}
uniformly in $\eps$ and in the finite fast sum.

Consider next the infinite product
\[
 \Theta_\eps(s)=\prod_{k\ge1}\frac{\lambda_k-s}{\lambda_k+s}.
\]
For each fixed $\eps$, \eqref{eq:sum} gives locally uniform convergence
both in the right half-plane and on compact boundary intervals.
Indeed, the deviation of a large-index factor from one is
$-2s/(\lambda_k+s)$, uniformly bounded by $C_K/\lambda_k$ on a
compact set $K$. The product is bounded by one and has unit boundary
modulus; hence it is inner. Every $\lambda_k$ remains a zero.

The brackets \eqref{eq:brackets} also give the sharper information
\begin{equation}\label{eq:sum-limit}
 D-8\eps\le\sum_{k\ge1}\lambda_k^{-1}\le D,\qquad
 \sum_{k\ge1}\lambda_k^{-3}\le16D\eps^2.
\end{equation}
To verify the first bound, compare the sums of
$h_\eps(r)=4\eps/[1+(2\pi\eps r/D)^2]$ at $k$ and $k+1$.
Its integral is $D$, the sum at $k\ge1$ differs from that integral
by at most $4\eps$, and omitting its first term loses at most another
$4\eps$. The second bound follows from $\lambda_k\ge1/(4\eps)$.
For every fixed $y\in\mathbb R$,
\begin{equation}\label{eq:Theta-limit}
 \Theta_\eps(iy)
 =\exp\!\left(-2i\sum_{k\ge1}\arctan(y/\lambda_k)\right)
 \longrightarrow e^{-2iDy}.
\end{equation}
The phase error after replacing each arctangent by its argument is
at most $(2/3)|y|^3\sum\lambda_k^{-3}$, since
$|\arctan r-r|\le|r|^3/3$.
Both multipliers have modulus one on the boundary.
Dominated convergence therefore gives strong convergence of their
boundary multiplication operators, and of their conjugates.
Formula~\eqref{eq:projection} now proves
\begin{equation}\label{eq:projection-limit}
 P_{\Theta_\eps}\longrightarrow P_{e^{-2Ds}}
       \quad\hbox{strongly on }\HH.
\end{equation}

Under inverse Laplace transformation, multiplication by $e^{-2Ds}$
is the right shift by $2D$. Its model space is exactly the transform
of $L^2(0,2D)$.
In the following time-domain statement the projection means
$\LL^{-1}P_\Theta\LL$.
Every finite fast observation belongs to $K_{\Theta_\eps}$,
since its transform is a combination of the kernels at the product's
zeros.
Consequently, any weak $L^2(0,\infty)$ limit of a bounded sequence of
finite fast observations as $\eps\downarrow0$ is supported in $[0,2D]$.
For a fixed test $g$ vanishing on $(0,2D)$,
\[
 \ip{f_\eps}g=\ip{f_\eps}{P_{\Theta_\eps}g}\longrightarrow0
\]
by \eqref{eq:projection-limit} and the uniform bound on $f_\eps$.
This proves the support assertion without requiring norm convergence
of the observations.

\subsection{Separating the slow mode and closing the upper bound}

\begin{lemma}\label{lem:distance}
For every $T>2D$ there are $\delta_{L,T}>0$ and $\eps_0>0$ such that
\begin{equation}\label{eq:distance}
 \inf_{f\in\Span\{e^{-\lambda_kt}:k\ge1\}}
       \norm{e^{-\lambda_0t}-f}_{(0,T)}\ge\delta_{L,T},
       \qquad 0<\eps<\eps_0,
\end{equation}
where the span consists of finite linear combinations.
\end{lemma}
\begin{proof}
If no such uniform distance existed, one could choose $\eps_j\to0$
and finite fast observations $f_j$ converging to
$e^{-\lambda_0(\eps_j)t}$ in $L^2(0,T)$.
Since $\lambda_0(\eps_j)\to0$, they converge to one there.
Estimate~\eqref{eq:comparison} bounds them in $L^2(0,\infty)$.
A weakly convergent subsequence exists in this Hilbert space
\cite[Lemma 1.13]{Teschl}.
Its limit is supported in $[0,2D]$ by
\eqref{eq:projection-limit}, but equals one on $(0,T)$.
This is impossible because $T>2D$.
\end{proof}

A full finite adjoint sum has flux
$q(t)=\alpha e^{-\lambda_0t}+f(t)$, with $f$ in the fast span.
Equation~\eqref{eq:distance} and $\lambda_0>0$ give
\[
 |\alpha|\le\delta_{L,T}^{-1}\norm q_{(0,T)},\qquad
 \norm f_{(0,T)}\le(1+\sqrt T/\delta_{L,T})\norm q_{(0,T)}.
\]
For $T>2\sqrt2D$, the mode bounds, \eqref{eq:fast-upper} and
\eqref{eq:comparison} thus imply
\begin{equation}\label{eq:observation}
 \norm{z(T)}_2\le
 \left(\frac{2\sqrt D}{\delta_{L,T}}+
       C_{L,T}\eps^{-1/2}e^{-(T-2\sqrt2D)/(4\eps)}\right)
                    \norm q_{(0,T)}.
\end{equation}
This estimate is uniform in the truncation and in all sufficiently small
$\eps$. At each fixed $\eps$, the graph-norm approximation and trace
convergence in Section~\ref{sec:density} extend it to every
$z_0\in X_\eps$. Appendix~\ref{sec:duality} then supplies a real null
control for every $L^2$ initial datum, with a uniform norm bound.
Since $2\sqrt2D=4\sqrt2L$, this proves the first assertion of
Theorem~\ref{thm:main}.

\section{A matching lower bound}\label{sec:lower}

\subsection{An inverse-Gram column with coherent spatial signs}

Take $N=\lfloor\eps^{-2}\rfloor$ and the fast Gram matrix of
\eqref{eq:Gram}. Choose the real coefficient vector $\boldsymbol c=G^{-1}e_1$,
where $e_1$ is the first coordinate vector. Then
\begin{equation}\label{eq:column}
 c_k=\frac{4(-1)^{k+1}\lambda_k\lambda_1}
 {(\lambda_k+\lambda_1)d_{k,N}d_{1,N}},\qquad
 c_1=\frac{2\lambda_1}{d_{1,N}^2}>0.
\end{equation}
For the smooth adjoint test
$z(t)=\sum_{k=1}^Nc_ke^{-\lambda_kt}\chi_k$,
left normalization gives $q(t)=\sum c_ke^{-\lambda_kt}$ and
\begin{equation}\label{eq:column-norm}
 \norm q_{(0,\infty)}^2=\boldsymbol c^TG\boldsymbol c=c_1.
\end{equation}

We need a single spatial layer in which all the terms contributing to
$z(T)$ have one sign. Put $x=L-s$ and normalize
$S_{\lambda_k}(L-s)$ by $-S_{\lambda_k}'(L)$.
The resulting function $y_k$ solves
\[
 y_k''+Q_k(s)y_k=0,\qquad y_k(0)=0,\quad y_k'(0)=1,\qquad
 Q_k(s)=\xi_k^2+\frac{\sech^2((L-s)/(2\eps))}{2\eps^2}.
\]
By \eqref{eq:brackets}, for all $k\le N$ and small $\eps$,
$0\le Q_k\le M_L^2\eps^{-4}$.
Put $\iota(s)=s$. The Volterra equation is $y_k=\iota-K_ky_k$, where
$(K_kg)(s)=\int_0^s(s-r)Q_k(r)g(r)\,dr$.
Writing $m=M_L/\eps^2$, induction gives
\[
 |(K_k^n\iota)(s)|\le\frac{m^{2n}s^{2n+1}}{(2n+1)!},\qquad n\ge0.
\]
The convergent Volterra series
$y_k=\sum_{n\ge0}(-1)^nK_k^n\iota$ therefore yields
\[
 |y_k(s)-s|\le s\left(\frac{\sinh(ms)}{ms}-1\right).
\]
Choose a fixed $\rho_L>0$ so that
$\sinh(M_L\rho_L)/(M_L\rho_L)-1\le1/2$.
For $0<s<\rho_L\eps^2$ we have $y_k(s)\ge s/2$
simultaneously for every $k\le N$.
In this layer
$c(L-s)/c(L)\ge e^{-s/(2\eps)}\ge1/2$ for small $\eps$.
The exact derivative signs in \eqref{eq:endpoint} now give
\begin{equation}\label{eq:layer}
 \operatorname{sgn}\chi_k(L-s)=(-1)^k,\qquad
 |\chi_k(L-s)|\ge\frac{s}{4\eps},
 \quad 1\le k\le N,\quad 0<s<\rho_L\eps^2.
\end{equation}
All products $c_k\chi_k(L-s)$ are negative.
Keeping only the first term when bounding their sum from below gives
\begin{equation}\label{eq:spatial-lower}
 \norm{z(T)}_2^2
 \ge\frac{\rho_L^3}{48}\eps^4c_1^2e^{-2\lambda_1T}.
\end{equation}
The common layer and the exact signs prevent cancellation; no
orthogonality of the adjoint eigenfunctions is used.

\subsection{The scalar product and the cost exponent}

Set $a_0=\pi^2/D^2$ and $R_\eps^2=1+1/(2a_0\eps^2)$.
For $v>u>0$, $(v-u)/(v+u)$ increases with $v$ and decreases with $u$.
Applying \eqref{eq:brackets} to each factor gives
\begin{equation}\label{eq:scalar-product}
 d_{1,N}\le\prod_{j=3}^{N+1}\frac{j^2-1}{j^2+R_\eps^2}.
\end{equation}
The corresponding infinite product is
\begin{equation}\label{eq:sinh-product}
 \prod_{j=3}^\infty\frac{j^2-1}{j^2+R_\eps^2}
 =\frac23\,\frac{\pi R_\eps(1+R_\eps^2)(1+R_\eps^2/4)}
                              {\sinh(\pi R_\eps)}.
\end{equation}
Indeed, $\prod_{j=3}^M(1-1/j^2)=2(M+1)/(3M)\to2/3$,
while $\sinh z=z\prod_{j\ge1}(1+z^2/(\pi^2j^2))$
\cite[Equation 4.36.1]{DLMF} gives the denominator after removing
its first two factors.
The ratio between the finite product in \eqref{eq:scalar-product}
and \eqref{eq:sinh-product} is at most
\[
 \exp\!\left((R_\eps^2+1)
             \sum_{j>N+1}\frac1{j^2-1}\right)\le C_L,
\]
because $N\ge1/(2\eps^2)$ and the tail sum is at most $1/N$.
Since $\pi R_\eps\ge D/(\sqrt2\eps)$ and $R_\eps\le C_L/\eps$,
\begin{equation}\label{eq:d1-lower}
 d_{1,N}\le C_L\eps^{-5}e^{-D/(\sqrt2\eps)}.
\end{equation}
Thus the truncation has a uniformly controlled omitted tail.

Combining \eqref{eq:column-norm} and \eqref{eq:spatial-lower} gives
\[
 \frac{\norm{z(T)}_2^2}{\norm q_{(0,T)}^2}
 \ge C_L\eps^4c_1e^{-2\lambda_1T}
 =C_L\eps^4\frac{2\lambda_1}{d_{1,N}^2}e^{-2\lambda_1T}.
\]
The brackets imply
$1/(4\eps)\le\lambda_1\le1/(4\eps)+4a_0\eps$.
Using \eqref{eq:d1-lower} therefore yields
\begin{equation}\label{eq:lower-ratio}
 \frac{\norm{z(T)}_2^2}{\norm q_{(0,T)}^2}
 \ge c_{L,T}\eps^{13}
    \exp\!\left(\frac{2D/\sqrt2-T/2}{\eps}\right).
\end{equation}
The cost inequality \eqref{eq:ratio} proves \eqref{eq:lower-main}.
Its right side diverges at every $0<T<2\sqrt2D=4\sqrt2L$.
Together with \eqref{eq:observation}, this proves
$T_{\rm unif}=4\sqrt2L$ and completes Theorem~\ref{thm:main}.

\appendix
\section{Transposition solutions and exact null controls}\label{sec:duality}

We give the fixed-viscosity functional-analytic details to retain the
full $L^2$ initial-data and boundary-control classes in \eqref{eq:cost}.
The dual $X_\eps'$ is the real continuous dual, or the continuous
anti-dual after complexification. The graph norm on $X_\eps$ is
$\norm v_2+\norm{A_\eps^*v}_2$.

For smooth data and a boundary control compactly supported in $(0,T)$,
integration by parts gives \eqref{eq:duality}.
The right boundary term vanishes, and the left term is
$+\eps h(t)\overline{z_x(T-t,-L)}$.
The spectral semigroup and trace estimate of
Section~\ref{sec:density} show, at fixed $\eps$, that
$\sup_{0\le t\le T}|q(t)|\le C_{\eps,L,T}\norm{z_0}_{X_\eps}$.
Consequently \eqref{eq:duality} defines a bounded terminal functional
in $X_\eps'$ for all $u_0\in L^2$ and $h\in L^2(0,T)$.
The same formula at each intermediate time gives a family in
$C([0,T];X_\eps')$.
Smooth approximations converge uniformly in this space.
Classical approximating solutions are obtained by lifting the boundary
with $h(t)\ell(x)$, where $\ell(-L)=1$, $\ell(L)=0$, and solving the
homogeneous-boundary forced equation by the spectral semigroup.
The term containing $h'$ is removed by time integration by parts
before taking the $L^2$ limit.

This solution also lies in spatial-temporal $L^2$.
For $F\in L^2((0,T)\times(-L,L))$ solve the backward test problem
\[
 -v_t+A_\eps^*v=F,\qquad v(T)=0,\qquad v(t,\pm L)=0.
\]
After the similarity by $c$, the orthonormal coefficients satisfy
$v_j(t)=\int_t^T e^{-\lambda_j(s-t)}F_j(s)\,ds$.
Weighted Cauchy--Schwarz and Tonelli give
\[
 \int_0^T|\lambda_jv_j(t)|^2\,dt
 \le\int_0^T|F_j(t)|^2\,dt.
\]
For example, use the nonnegative weight
$\lambda_je^{-\lambda_j(s-t)}$, whose integral in either variable
is at most one. Parseval, the graph/$H^2$ equivalence, and the
similarity then imply
\[
 \norm v_{L^2(0,T;H^2)}+\norm{v(0)}_2+
     \norm{\eps v_x(\cdot,-L)}_{(0,T)}
 \le C_{\eps,L,T}\norm F_{L^2((0,T)\times(-L,L))}.
\]
The $L^2$ bound on $v$ and its initial value follows also from
the contraction of the transformed semigroup, giving respectively
$T\norm F$ and $\sqrt T\norm F$ before the similarity.
Riesz representation applied to
\[
 F\longmapsto
 \ip{u_0}{v(0)}+\int_0^T h(t)\overline{\eps v_x(t,-L)}\,dt
\]
constructs a unique spacetime $L^2$ solution.
For smooth data it is the lifted classical solution.
The estimate just proved and the earlier terminal estimate show that
the spacetime and $C([0,T];X_\eps')$ limits agree.
They preserve the original initial datum and both boundary values in
transposition. Zero terminal functional is therefore exactly the zero
terminal distribution, and hence the zero terminal state.
The standard Sobolev and Hilbert-space facts used here are recorded,
for example, in \cite[Sections 0.3, 1.1--1.3 and 2.7]{Teschl}.

Finally suppose
\[
 \norm{z(T)}_2\le M\norm q_{(0,T)}\qquad(z_0\in X_\eps).
\]
On the range of observed fluxes define the linear functional
$\ell(q)=-\ip{z(T)}{u_0}$.
It is well defined by this inequality and has norm at most
$M\norm{u_0}_2$.
Extend it by continuity to the range closure, then compose with the
orthogonal projection onto that closed subspace.
Riesz representation gives $g\in L^2(0,T)$ with
$\ell(q)=\ip qg$ and $\norm g_{(0,T)}\le M\norm{u_0}_2$.
Setting $h(t)=g(T-t)$ and conjugating this equality in
\eqref{eq:duality} annihilates every terminal test.
For real $u_0$, the observation range and functional respect complex
conjugation; uniqueness of the Riesz representative makes $g$ real.
No closedness of the original observation range is assumed.
Conversely, for a real smooth adjoint test with $z(T)\ne0$, choose
$u_0=z(T)/\norm{z(T)}_2$ in \eqref{eq:duality}.
Every null control satisfies
$\norm h_{(0,T)}\norm q_{(0,T)}\ge\norm{z(T)}_2$,
which proves \eqref{eq:ratio}.

For completeness, the exact scaling
$y=x/L$, $s=t/L$, $\delta=\eps/L$,
$v(s,y)=\sqrt L\,u(Ls,Ly)$ and $H(s)=\sqrt L\,h(Ls)$
preserves both initial and control norms.
Thus, writing the interval dependence explicitly,
$C_L(T,\eps)=C_1(T/L,\eps/L)$.
The proof above treats arbitrary fixed $L$ directly.

\section*{Declaration of AI assistance}
AI assistance was used in literature retrieval, developing and checking
the argument, and preparing the manuscript. The mathematical claims
and proofs are presented here for independent review.

\begin{thebibliography}{9}
\bibitem{Laheurte}
V. Laheurte,
\emph{Cost of controllability of the Burgers equation linearized at a
steady shock in the vanishing viscosity limit},
SIAM Journal on Mathematical Analysis \textbf{58} (2026), 2702--2737.
\href{https://doi.org/10.1137/25M1723311}{doi:10.1137/25M1723311};
\href{https://arxiv.org/abs/2411.12267v3}{arXiv:2411.12267v3}.

\bibitem{KoikeLaheurte}
K. Koike and V. Laheurte,
\emph{Uniform null-controllability times for the Coron--Guerrero
problems are $2$ and $2+2\sqrt2$},
preprint (2026).
\href{https://arxiv.org/abs/2609.35355v1}{arXiv:2609.35355v1}.

\bibitem{BaranovEtAl}
A. Baranov, P. Jaming, K. Kellay and M. Speckbacher,
\emph{Oversampling and Donoho--Logan type theorems in model spaces},
Annales Fennici Mathematici \textbf{49} (2024), 167--182.
\url{https://afm.journal.fi/article/view/143957}.

\bibitem{Teschl}
G. Teschl,
\emph{Mathematical Methods in Quantum Mechanics:
With Applications to Schr\"odinger Operators},
second edition, Graduate Studies in Mathematics 157,
American Mathematical Society, 2014.
\url{https://www.mat.univie.ac.at/~gerald/ftp/book-schroe/schroe2.pdf}.

\bibitem{Weikard}
R. Weikard, \emph{Complex Analysis}, MA 648 lecture notes,
University of Alabama at Birmingham, December 8, 2016.
\url{https://www.uab.edu/cas/mathematics/images/Documents/weikard-complex-analysis.pdf}.

\bibitem{DLMF}
\emph{NIST Digital Library of Mathematical Functions},
version 1.2.8, September 15, 2026.
Sections 1.10(ix) and 4.36.
\url{https://dlmf.nist.gov/}.
\end{thebibliography}
\end{document}
